% Generated by code/rebuild/WriteResults.py from corrected-run outputs.
\section{Results}

We first asked whether a lasting change in how some traders respond to
their neighbors could turn a mixed market into one dominated by sellers.
Rule~19286 differs from the baseline Rule~19280 in a single response:
a trader with three neighbors in Buy and one in Hold chooses Sell rather
than Buy (Fig.~\ref{fig:rule-quench}A,B). We paired each switched market
with a no-switch control that kept Rule~19280 and shared its baseline
history (Supplementary
Table~\ref{tab:campaigns}). With active technical traders and dealer
depth $h=100$, none of the 50 starting markets reached a Sell majority
when the fraction of interior traders given the new rule was $f=0.10$ or
0.15, whereas all of them did at $f=0.20$ and 0.25. No control
reached a majority over the same period. The median onset fell from
1,070.0 steps at fraction $f=0.20$ to 311.5 steps at $f=0.25$
(Supplementary Table~\ref{tab:paired-switch}). In the displayed pair,
groups of sellers expanded after the rule switch while the control kept
its mixture of decisions (Fig.~\ref{fig:rule-quench}C). A one-time
instruction did not have the same effect: setting the selected traders
to Sell once at fraction $f=0.15$, 0.20, or 0.25 produced no majority in either technical-trader condition
(Supplementary Table~\ref{tab:revision-pulse}). Which response changed
also mattered. Of the eighteen alternative rules tested at fraction
$f=0.20$, each differing from Rule~19280 in a single response,
6 produced a majority in every tested starting market under both technical-trader conditions, while 12 produced none (Supplementary Table~\ref{tab:revision-rules}). Thus,
some tested lasting rule switches produced majorities where one-time Sell
instructions did not, and outcomes depended on which response changed
and how many traders received it.

Because technical traders respond to the quoted price, a Sell majority
might have required price feedback. To determine whether it did, we
repeated the rule switch with every technical trader held in Hold, so
that trades still completed and the price still moved but no decision
responded to it (Fig.~\ref{fig:reservoir}).
With technical traders held, the Sell-majority probability was
0.34 at fraction $f=0.175$, 0.70 at $f=0.185$, and
1.00 at $f=0.20$ (Fig.~\ref{fig:reservoir}A). For each
starting market and fraction, onset, final mean Sell fraction, and final
interior grid were identical across dealer depths despite different price
paths. Without technical-trader responses, financial variables cannot influence decisions,
so responses to neighbors were sufficient to produce a Sell majority.
Allowing technical traders to respond to price changed majority probability
near the crossover. At dealer depth $h=100$, 25 of 50 markets reached a majority
with technical traders active at fraction $f=0.175$, compared with 17
when held. At fraction $f=0.185$, the corresponding counts were 42 and 35 (Fig.~\ref{fig:reservoir}A,B).
Both the size and the sign of the difference in majority probability between active and held conditions depended
on the dealer depth and fraction. Thus, price feedback could change how
often the rule switch produced a majority, but it was not needed for
selling decisions to spread.

\input{sections/finite-size-results}

A 3,000-step observation period could end before a slow majority formed,
so we extended five markets with active technical traders and dealer
depth $h=100$ to 10,000 steps (Fig.~\ref{fig:dose-panels}). Selection used the
3,000-step outcomes, before image inspection. Two had not reached a
majority by that endpoint and three had. Of the two without a majority,
the market at fraction $f=0.175$ reached one after 4{,}073 steps, whereas
the market at $f=0.150$ had not reached one by 10,000 steps. The three
earlier majorities formed after 2{,}330, 1{,}836, and 1{,}028 steps
at fractions $f=0.175$, 0.185, and 0.20, respectively (Supplementary Table~\ref{tab:market-panel-horizons}). Each
displayed majority formed as groups of sellers grew and merged (Fig.~\ref{fig:dose-panels}A,C). For the market that reached a majority late,
maximum drawdown rose from 10.17\% at 3,000 steps to 26.34\% at
10,000 steps (Supplementary Table~\ref{tab:market-panel-horizons}). To measure late-majority
frequency at fraction
$f=0.150$, we extended all 50 starting markets
and their no-switch controls. Two switched markets
reached a majority by 10,000 steps, compared with none at 3,000 steps,
and no control reached one. Among the switched markets, the median
drawdown changed from 5.04\% to 5.87\%, whereas the largest
drawdown rose from 12.29\% to 25.30\%
(Supplementary Table~\ref{tab:market-panel-horizons}). Thus, remaining mixed
for 3,000 steps did not establish a long-run state; majority outcomes
depended on the observation period.

We next asked how dealer reserves affected majority probability and
the resulting price decline. We varied the dealer depth $h$, the dealer's initial shares per
trader, while keeping the same rule switch (Fig.~\ref{fig:reservoir}). Near the crossover, dealer depth also affected majority probability.
We tested whether the technical traders' adaptive band widths contributed
to that dependence. At fraction $f=0.175$, the Sell-majority
probability was 0.44 at dealer depth $h=0.1$ and 0.50 at $h=100$, a
difference of -0.06 that became -0.26 when the band widths were frozen
after the rule switch. We matched starting markets across depths and band conditions.
Subtracting the frozen-width contrast from the adaptive-width contrast
gave a paired difference of 0.20 with a 95\% bootstrap interval of
[-0.08, 0.50], leaving the contribution of band-width adaptation unresolved (Supplementary Table~\ref{tab:revision-band}). At fraction
$f=0.20$, every starting market reached a majority at every tested
depth, yet with active technical traders the median maximum drawdown
was 99.99\% at dealer depth $h=0.1$, 99.52\% at $h=1$, and 25.15\% at $h=100$
(Fig.~\ref{fig:reservoir}C). In an additional comparison, median drawdowns were
48.62\% at dealer depth $h=40$ and 7.08\% at $h=400$
(Supplementary Table~\ref{tab:revision-impact}). Smaller reserves made the quoted price
more sensitive to net selling, because the reserves set both the
dealer's available funds and the price required to execute a net trade.
Much of the decline at dealer depth $h=0.1$ occurred without the rule
switch: no-switch controls there already had a median drawdown of
77.04\%, compared with 0.75\% at $h=100$. The median paired
increase in drawdown relative to the no-switch control was 22.96 percentage points at dealer depth
$h=0.1$ and 24.39 at $h=100$ (Supplementary Table~\ref{tab:paired-switch}).
Small positive reserves accompanied these declines while rejecting few
orders backed by trader resources, with median dealer-rejection fractions of 0.000226 at dealer
depth $h=0.1$, 0.000034 at $h=1$, and 0.0000006 at $h=100$
(Fig.~\ref{fig:reservoir}D). Zero reserves were a separate limit, in which only
opposing trader orders matched, the dealer's net trading was zero, and
the price stayed fixed. In the
displayed starting market 2, a majority formed at dealer depths
$h=0.1$, 1, and 100 with active technical traders while the price followed distinct
paths (Fig.~\ref{fig:dealer-panels}). Thus, across the tested depths at fraction $f=0.20$, the dealer's
reserves set the severity of the decline rather than whether a majority
formed, and shallow positive reserves did not prevent severe drawdowns.

Trades also depended on the traders' own resources, so we independently
changed their cash and shares at the rule switch (Fig.~\ref{fig:balance-sheet}). With technical traders held, every cash and
share combination at fraction $f=0.18$ produced the same decision
history within a starting market, and the Sell-majority probability was
0.64 at every dealer depth. Holdings entered the order and clearing
calculations but not the interior decision rule, so they could change
decisions only through the price response of active technical traders.
With technical traders active, the probability ranged from
0.54 to 0.80 across resource combinations and dealer depths
(Fig.~\ref{fig:balance-sheet}A; Supplementary Table~\ref{tab:balance-fractions}). With unchanged
cash and dealer depth $h=100$, share multipliers of 0.1, 1, and 10 gave
median drawdowns of 3.17\%, 20.18\%, and 77.55\%
(Fig.~\ref{fig:balance-sheet}B), because more shares supplied more potential sales
while the reserve curve set the price at which net selling could clear.
By contrast, raising the cash multiplier from 1 to 10 at dealer depth
$h=1$ with unchanged shares changed the median drawdown only from
96.46\% to 97.10\% and the majority probability from
0.58 to 0.56, in markets where the decline was already close to
complete (Fig.~\ref{fig:balance-sheet}A,B). Few dealer rejections did not imply a small drawdown
across these resource combinations (Fig.~\ref{fig:balance-sheet}C,D). Together,
these comparisons showed that financial outcomes in switched markets
depended on traders' cash and shares and on the dealer's reserves. A majority decision
to sell did not by itself determine how many shares could change hands
or how far the quoted price would fall.
